\documentclass[main]{subfiles}

\title[Advanced Hyperparameter Optimization]{Advanced Hyperparameter Optimization}
\subtitle{Warmstarting via Meta-Learning Across Tasks}
\author[Katharina Eggensperger]{Katharina Eggensperger}
\date{}

\titlebackground*{images/AutoML_Speeding_up_HPO.jpg}

%\institute{}
%\date{}

\begin{document}
	
	\maketitle
	
%----------------------------------------------------------------------
%----------------------------------------------------------------------
\begin{frame}[c]{What if we have more tasks?}

\begin{columns}[T,onlytextwidth]
    \column{0.50\textwidth}
    \vspace{10pt}
    In practice, we often run HPO on similar but slightly different datasets, e.g., 
    \begin{itemize}
        \item to track changes in data(streams),
        \item to adapt to similar tasks.
    \end{itemize}
    The best setting tends to remain stable across tasks.
    \column{0.50\textwidth}
    \vspace{-1cm}
    \begin{figure}
        \centering
        \includegraphics[width=\linewidth]{05_Advanced_HPO/images/meta-learning_hpo.png}
        \caption{Adapted from: \liturl{Schneider et al. PPSN'2022}{https://arxiv.org/abs/2208.00220}}
    \end{figure}
\end{columns}

\vspace{0.5cm}
\textbf{$\leadsto$ How can we increase the efficiency of HPO on a sequence of \\related tasks?} \\

\end{frame}
%-----------------------------------------------------------------------%-----------------------------------------------------------------------%----------------------------------------------------------------------
\begin{frame}[c]{A Motivating Example}
    \begin{figure}
        \centering
        \includegraphics[width=0.4\linewidth]{05_Advanced_HPO/images/HPO_multidata.png}
        \caption{Adapted from: \liturl{Greff et al. 2017}{https://arxiv.org/abs/1503.04069}}
    \end{figure}
\end{frame}
%-----------------------------------------------------------------------
%----------------------------------------------------------------------

\begin{frame}[c]{Meta-Learning Across Tasks}

Assuming a fixed search space (i.e., HPO task), a given target dataset $\dataset_t$, and a given set of prior datasets $\datasets_{meta} = \{\dataset_1, \dots, \dataset_{t-1}\}$ on which we have collected prior observations.

\vspace{0.5cm}
\textbf{Why?} Objective functions of tasks are correlated. \\
\textbf{Idea:} Positive transfer from similar tasks, zero transfer from different tasks. \\
\textbf{Method:} Portfolios to (partially) replace the initial design.

\vspace{0.5cm}
\pause
\textbf{Option 1:} \textbf{Task-agnostic} $\leadsto$ well-performing regions of "all" tasks. \\
\textbf{Option 2:} \textbf{Task-specific} $\leadsto$ well-performing regions on similar tasks.

\end{frame}
%-----------------------------------------------------------------------
%----------------------------------------------------------------------
\begin{frame}{Task-Agnostic Portfolios}

Given meta-datasets $\datasets_{meta}$ and a candidate pool $\confs_{cand} = \{\conf_1, ..., \conf_n\}$. 

The goal is to construct a portfolio $\portfolio \subseteq \confs_{cand}, |\portfolio| = k$, that minimizes the objective:

\begin{equation*}
\loss(\portfolio, \datasets_{meta}) = \sum_{i=1}^{t-1} \min_{\conf \in \portfolio} \cost(\conf, \dataset_i)
\end{equation*}

\pause
\textbf{Key idea:} The portfolio
\begin{itemize}
    \item should include at least one well-performing configuration per task,
    \item covers optima across diverse tasks,
    \item and its performance is defined by the single best configuration per task.
\end{itemize}

\end{frame}
%-----------------------------------------------------------------------
%----------------------------------------------------------------------
\begin{frame}{Greedy Forward Selection}

\textbf{Intuition:}
\begin{itemize}
    \item Each step adds a configuration that improves performance on "uncovered" tasks
    \item Encourages diversity via complementary performance
\end{itemize}
\pause
\vspace{0.5cm}

\begin{columns}[T,onlytextwidth]
\begin{column}{0.45\textwidth}
\textbf{Input:} $\confs_{cand} = \{\conf_1, \dots, \conf_n\}$, $\datasets_{meta}$, $\cost(\cdot)$ \\
\textbf{Initialize:} $\portfolio \gets \emptyset$ \\
\begin{enumerate}
    \item \textbf{while $|\portfolio| < k$ do} \\
    \item \quad $\conf^+ \in \argmin_{\lambda \in \confs_{cand}} \loss \left(\portfolio \cup \conf, \datasets_{meta} \right)$ \\
    \item \quad $\portfolio \gets \portfolio \cup \conf^+, \confs_{cand} \gets \confs_{cand} \setminus \conf^+$ \\
    \item \textbf{return $\portfolio$}
\end{enumerate}
\end{column}
\pause
\begin{column}{0.45\textwidth}
{%
\begin{figure}
    \centering
    \includegraphics[width=0.6\linewidth]{05_Advanced_HPO/images/portfolio.png}
\end{figure}
}%
\end{column}
\begin{column}{0.2\textwidth}
    
\end{column}
\end{columns}
\vspace{0.5cm}

\end{frame}
%-----------------------------------------------------------------------
%----------------------------------------------------------------------
\begin{frame}{Task-Specific Portfolios}

\begin{columns}[T,onlytextwidth]
  \column{0.4\textwidth}
  \vspace{0pt}
  Instead of finding a global portfolio, we want a portfolio specifically for our target task $t$.

  \vspace{0.2cm}
  \textbf{Key idea} on how to adapt greedy portfolio construction:
  \begin{itemize}
      \item Find the $k$-closest tasks
      \item Choose best performing configurations on these tasks to warmstart optimization
  \end{itemize}
  \column{0.35\textwidth}
  \vspace{-2cm}
  \begin{figure}
      \centering
      \includegraphics[width=\linewidth]{05_Advanced_HPO/images/news_task_similarity.png}
      \caption{Adapted from \liturl{Schneider et al. PPSN'2022}{https://arxiv.org/abs/2208.00220}}
      \label{fig:placeholder}
  \end{figure}
  \column{0.05\textwidth}
\end{columns}

\pause
\vspace{-0.5cm}
\textbf{How do we define task similarity?} \\
(Intuition: If tasks are similar, similar configurations achieve similar \\ performance, i.e., if the datasets are similar, good configurations transfer.)

\end{frame}
%-----------------------------------------------------------------------
%----------------------------------------------------------------------
\begin{frame}{Meta-features as Dataset Descriptors}

Meta-features $\mathbf{z} = \{\phi_1(\dataset) \dots, \phi_d(\dataset)\}$ representing dataset characteristics, e.g.,
\begin{itemize}
    \item \textbf{simple} features (e.g., \#samples, \#features),
    \item \textbf{information-theoretic} meta-features (e.g., entropy of features and labels),
    \item \textbf{model-based} features (e.g. depth of a tree),
    \item and \textbf{landmarking} features (e.g., performance of simple, fast ML models)
\end{itemize}
can be used to describe datasets.

\vspace{0.5cm}
Task similarity can then be computed as $\text{sim}(t, t') = -\| \mathbf{z}_t - \mathbf{z}_{t'}\|$
\end{frame}
%-----------------------------------------------------------------------
%----------------------------------------------------------------------
\begin{frame}{Meta-feature-based Portfolios}

Now, we can build task-specific portfolios:
\begin{enumerate}
    \item Compute $\{\mathbf{z}_1, \cdots, \mathbf{z}_t\}$
    \item Compute $k$-closest datasets 
    \item Use the best performing configurations of the closest $k$ datasets as the portfolio
\end{enumerate}
   
\end{frame}
%----------------------------------------------------------------------
%----------------------------------------------------------------------
\begin{frame}{Task-specific vs Task-agnostic}

\setbeamercolor{block title}{use=normal text,fg=normal text.fg}%
\setbeamercolor{block body}{use=normal text,fg=normal text.fg}%
\begin{columns}[T,onlytextwidth]
  \column{0.45\textwidth}
  \vspace{0pt}
  \begin{block}{\raggedright Task-Agnostic}
\begin{itemize}
    \item can be built offline once
    \item covers diverse settings
    \item emphasizes complementarity in configurations
\end{itemize}
\end{block}
\pause
\column{0.45\textwidth}
\begin{block}{\raggedright Task-Dependent}
\begin{itemize}
    \item build per task
    \item requires hand-designed meta-features
    \item emphasizes performance
    \item configurations might be very similar
\end{itemize}
\end{block}
   \end{columns}

\end{frame}
%----------------------------------------------------------------------
%----------------------------------------------------------------------
\begin{frame}{Alternatives and Extensions}

Instead of meta-features, we can also use a
\begin{itemize}
    \item \textbf{ranking-based} similarity via comparing performance across shared configurations,
    \item or \textbf{learned} similarity via task embeddings, e.g, Dataset2Vec \liturl{Jomaa et al. 2021}{https://dl.acm.org/doi/abs/10.1007/s10618-021-00737-9} or embeddings of textual descriptions.
\end{itemize}

\pause
\vspace{0.5cm}
Instead of portfolios, we can also:
\begin{itemize}
    \item combine/ensemble \textbf{surrogate models} across tasks
    \item combine/ensemble \textbf{acquisition values} across tasks
\end{itemize}
\end{frame}
%----------------------------------------------------------------------
%----------------------------------------------------------------------
\begin{frame}{Conclusion and Key Takeaway}

Prior observations on \emph{related tasks} can increase the efficiency of Bayesian optimization, e.g., by warmstarting with a
\begin{itemize}
    \item \textbf{task-agnostic} portfolio covering diverse settings,
    \item \textbf{task-specific} portfolio consisting of configurations that performed well in similar tasks.
\end{itemize}
Task similarity can be defined using meta-features, configuration rankings, or learned embeddings.
\pause
\vspace{0.5cm}

$\leadsto$ If the assumed task similarity is incorrect, transfer can be harmful (negative transfer).

\end{frame}

\end{document}